% Standalone problem document, generated from the adjacent README.md.
% Compile with XeLaTeX. Mathematical content is maintained in Markdown.
\documentclass[11pt,a4paper]{article}
\usepackage[fontsize=10.5pt]{scrextend}
\usepackage[margin=22mm,headheight=15pt,headsep=9mm,footskip=11mm]{geometry}
\usepackage{fontspec}
\setmainfont{texgyrepagella}[Extension=.otf,UprightFont=*-regular,BoldFont=*-bold,ItalicFont=*-italic,BoldItalicFont=*-bolditalic]
\setsansfont{texgyreheros}[Extension=.otf,UprightFont=*-regular,BoldFont=*-bold,ItalicFont=*-italic,BoldItalicFont=*-bolditalic]
\setmonofont{texgyrecursor}[Extension=.otf,UprightFont=*-regular,BoldFont=*-bold,ItalicFont=*-italic,BoldItalicFont=*-bolditalic,Scale=MatchLowercase]
\usepackage{amsmath,amssymb}
\usepackage{unicode-math}
\setmathfont{texgyrepagella-math.otf}
\usepackage{xcolor,microtype,enumitem,needspace,fancyhdr}
\usepackage{bookmark}
\definecolor{ink}{HTML}{17344B}
\definecolor{accent}{HTML}{197D80}
\definecolor{muted}{HTML}{596773}
\hypersetup{colorlinks=true,linkcolor=accent,urlcolor=accent,pdftitle={MI-21 - Freewan--Hayajneh
inequality for sums of weighted geometric
means},pdfauthor={Open Problems in Numerical Linear Algebra}}
\urlstyle{same}
\setlength{\parindent}{0pt}
\setlength{\parskip}{5pt plus 1pt minus 1pt}
\linespread{1.04}
\setlength{\emergencystretch}{3em}
\widowpenalty=10000
\clubpenalty=10000
\displaywidowpenalty=10000
\allowdisplaybreaks[1]
\setlist{nosep,leftmargin=1.5em}
\providecommand{\tightlist}{\setlength{\itemsep}{0pt}\setlength{\parskip}{0pt}}
\providecommand{\pandocbounded}[1]{#1}
\pagestyle{fancy}
\fancyhf{}
\fancyhead[L]{\sffamily\footnotesize\color{muted}OPEN PROBLEMS IN NUMERICAL LINEAR ALGEBRA}
\fancyhead[R]{\sffamily\bfseries\footnotesize\color{accent}MI-21}
\fancyfoot[L]{\parbox[b]{0.9\textwidth}{\sffamily\fontsize{8}{10}\selectfont\color{muted}Editorial ratings; a bounded search cannot certify that a problem remains unsolved.\\Literature check: 2026-09-12}}
\fancyfoot[R]{\sffamily\footnotesize\color{muted}\thepage}
\renewcommand{\headrulewidth}{0.3pt}
\renewcommand{\headrule}{\hbox to\headwidth{\color{accent}\leaders\hrule height \headrulewidth\hfill}}
\makeatletter
\renewcommand{\section}{\Needspace{5\baselineskip}\@startsection{section}{1}{0pt}{12pt plus 2pt minus 2pt}{4pt}{\sffamily\bfseries\large\color{ink}}}
\renewcommand{\subsection}{\Needspace{4\baselineskip}\@startsection{subsection}{2}{0pt}{9pt plus 1pt minus 1pt}{3pt}{\sffamily\bfseries\normalsize\color{ink}}}
\makeatother
\setcounter{secnumdepth}{0}
\begin{document}
{\sffamily\small\color{muted}Matrix inequalities and norms\par}
\vspace{3pt}
{\raggedright\hyphenpenalty=10000\exhyphenpenalty=10000\sffamily\bfseries\fontsize{21}{24}\selectfont\color{ink}Freewan--Hayajneh
inequality for sums of weighted geometric means\par}
\vspace{7pt}
{\sffamily\small\textbf{Difficulty:} challenging\quad\textbf{Importance:} interesting
to the community\par}
{\sffamily\small\bfseries\color{accent}Status: Lean verified\par}
\vspace{4pt}
{\color{accent}\hrule height 0.8pt}
\vspace{6pt}
\textbf{Rating rationale:} The coupled power and weight parameters make
extension beyond the proved cases challenging; positive definite
aggregation and matrix-function estimates give community relevance.

\subsection{Resolution --- 2026-09-11}\label{resolution-2026-09-11}

\textbf{Negative result by Matthew J. Colbrook}, Department of Applied
Mathematics and Theoretical Physics, University of Cambridge.
\textbf{Independent proof review: PASS.}

Two rational positive definite \(2\times2\) summands with \(s=t=1/2\),
\(r=2\) and aggregate matrices \(A=B=I\) violate the operator-norm
inequality for every \(p>0\). The left side has eigenvalue
\(1351000/1350907>1\), while the right side is one.

The exact target is resolved. The complete negative answer now has
\href{https://github.com/ajt60gaibb/OpenProblemsInNLA/blob/main/matrix-inequalities-and-norms/MI-21\#lean-proof-and-verification-evidence--2026-09-12}{Lean
proof and verification evidence}. The original statement and source
evidence are retained below; its former difficulty rating is historical.

\textbf{Primary manuscript:}
\href{https://github.com/ajt60gaibb/OpenProblemsInNLA/blob/main/matrix-inequalities-and-norms/MI-21/solution.pdf}{complete
proof PDF},
\href{https://github.com/ajt60gaibb/OpenProblemsInNLA/blob/main/matrix-inequalities-and-norms/MI-21/solution.tex}{standalone
TeX}, Theorem 1.1 and its proof;
\href{https://github.com/ajt60gaibb/OpenProblemsInNLA/blob/main/matrix-inequalities-and-norms/MI-21/solution.md}{authorship
and scope}. The
\href{https://github.com/ajt60gaibb/OpenProblemsInNLA/blob/main/references/colbrook-matrix-2026-09-11/verification/reviews/MI-21-review.md}{independent
review} checks the full original argument and records its hash. The
draft was AI-assisted; the 2026-09-11 review was independent agent
verification, without external human peer review or formal
certification. The subsequent Lean verification below covers the
complete negative answer.
\href{https://github.com/ajt60gaibb/OpenProblemsInNLA/blob/main/references/colbrook-matrix-2026-09-11/README.md}{Submission
record}.

\subsection{Lean proof and verification evidence ---
2026-09-12}\label{lean-proof-and-verification-evidence-2026-09-12}

\textbf{Mathematical counterexample:} Matthew J. Colbrook, Department of
Applied Mathematics and Theoretical Physics, University of Cambridge.
\textbf{Lean formalization:} George Stepaniants, Department of Computing
and Mathematical Sciences, California Institute of Technology, Pasadena,
California, USA, with AI-agent assistance.

The
\href{https://github.com/sgstepaniants/OpenProblemsInNLA/tree/06ade659dee260a18b79ce638383bf0a49125ecf/matrix-inequalities-and-norms/MI-21/lean}{proof
at revision 06ade65} establishes the full negative answer. It preserves
every original complex-matrix, dimension, parameter and unitarily
invariant norm quantifier. The witness uses the genuine Euclidean
operator norm, separately proved admissible, and actual spectral powers
and weighted geometric means. All input positivity and numerical
identities are proved. A nonzero eigenvector gives a strict norm
violation for every \(p>0\); it does not replace the norm definition.

\href{https://github.com/sgstepaniants/OpenProblemsInNLA/blob/06ade659dee260a18b79ce638383bf0a49125ecf/matrix-inequalities-and-norms/MI-21/lean/Solution.lean}{Solution.lean}
exports:

\begin{itemize}
\tightlist
\item
  \textit{NLA.MI21.operatorNorm\_isUnitaryInvariant}: the actual complex
  Euclidean operator norm satisfies every stipulated norm and two-sided
  unitary-invariance property.
\item
  \textit{NLA.MI21.counterexample}: all admissible witness hypotheses,
  genuine functional-calculus identities and the strict norm violation
  for every positive outer parameter.
\item
  \textit{NLA.MI21.not\_geometricMeanNormConjecture}: negation of the
  complete original universal assertion.
\end{itemize}

Two independent agents
\href{https://github.com/ajt60gaibb/OpenProblemsInNLA/blob/main/matrix-inequalities-and-norms/MI-21/lean/reviews}{reviewed
the frozen statements and complete proof} against the original target.
\href{https://github.com/sgstepaniants/OpenProblemsInNLA/actions/runs/34709291489}{Linux
run 34709291489} executed the actual sandboxed Lean4 Comparator on
GitHub Actions Ubuntu 24.04, matched all three declarations and replayed
the solution through Lean's default kernel. The
\href{https://github.com/ajt60gaibb/OpenProblemsInNLA/blob/main/matrix-inequalities-and-norms/MI-21/lean/verification/linux-2026-09-12}{original
artifacts and independent operational audit} retain the complete
81-input source identity, artifact digests and successful sandbox,
raw-kernel and axiom-rejection controls. All 11 internal/public
\href{https://github.com/ajt60gaibb/OpenProblemsInNLA/blob/main/matrix-inequalities-and-norms/MI-21/lean/verification/proof-axioms.json}{transitive
axiom reports} use only \textit{propext}, \textit{Classical.choice} and
\textit{Quot.sound}. This catalog reviewed the remote Linux execution;
local macOS re-elaborations are recorded separately. These are
independent AI-agent reviews, not external human peer review or
source-author endorsement. The statement referees establish
correspondence with the original prose; Comparator checks formal
identity and kernel trust.

The proof pins \textbf{Lean 4.33.1},
\href{https://github.com/alerad/leancert/tree/621a43d7cf21f87872392a01e874f2f1dbddc926}{LeanCert
621a43d} and
\href{https://github.com/leanprover-community/mathlib4/tree/0df444a360eaa60ab8c11dca51a86af692955474}{Mathlib
0df444a}, with
\href{https://github.com/ajt60gaibb/OpenProblemsInNLA/blob/main/matrix-inequalities-and-norms/MI-21/lean/lake-manifest.json}{all
dependency revisions locked}. A generic positive-square-root argument
identifies the actual geometric means from two exact Riccati
certificates. One kernel-mode LeanCert point certificate proves the
strict scalar gap; no interval subdivision or numerical matrix-root
approximation is needed. The
\href{https://github.com/ajt60gaibb/OpenProblemsInNLA/blob/main/matrix-inequalities-and-norms/MI-21/lean/formalization.yaml}{formalization
manifest},
\href{https://github.com/ajt60gaibb/OpenProblemsInNLA/blob/main/matrix-inequalities-and-norms/MI-21/lean/NUMERICAL_TARGETS.md}{numerical
obligations} and
\href{https://github.com/ajt60gaibb/OpenProblemsInNLA/blob/main/matrix-inequalities-and-norms/MI-21/lean/README.md}{project
guide} record scope and reproduction. The narrower parameter regimes
already proved in the literature are not contradicted.

From the verified proof revision, on an isolated non-root Linux host
meeting the
\href{https://github.com/ajt60gaibb/OpenProblemsInNLA/blob/main/tools/lean/HARNESS.md}{shared
harness prerequisites}, run:

\begin{verbatim}
tools/lean/bootstrap.sh /absolute/path/to/nla-lean-tools
tools/lean/selftest.sh /absolute/path/to/nla-lean-tools
tools/lean/verify.sh \
  matrix-inequalities-and-norms/MI-21/lean \
  /absolute/path/to/nla-lean-tools
\end{verbatim}

\subsection{Problem statement}\label{problem-statement}

For positive definite matrices, define

\[A\#_tB=A^{1/2}(A^{-1/2}BA^{-1/2})^tA^{1/2}.\]

Is the following true for all positive integers \(m,n\), all
\(A_i,B_i\in\mathbb C^{n\times n}\) positive definite, all
\(t\in[0,1]\), and all real \(s,r,p>0\) with \(sr\ge1\)? Set
\(A=\sum_iA_i\), \(B=\sum_iB_i\). For every unitarily invariant norm on
\(n\times n\) matrices, does

\[\left\|\sum_{i=1}^m(A_i^s\#_tB_i^s)^r\right\|
\le
\left\|\left(A^{(1-t)srp/2}B^{tsrp}A^{(1-t)srp/2}\right)^{1/p}\right\|\]

hold? Here all real powers of positive definite matrices use spectral
functional calculus, and unitary invariance means \(\|UXV\|=\|X\|\) for
unitary \(U,V\).

The conjecture controls the norm of a sum of nonlinear positive definite
matrix means using only sums of the input matrices. It is relevant to
aggregation and interpolation of positive definite data and to estimates
for algorithms involving matrix powers. The full weighted statement is
one problem; its individual parameter cases are not separate entries.

\subsection{References}\label{references}

\begin{enumerate}
\def\labelenumi{\arabic{enumi}.}
\tightlist
\item
  S. Freewan and M. Hayajneh, \emph{On a conjecture related to the
  geometric mean and norm inequalities}, Math. Inequal. Appl. 27 (2024),
  Conjecture 3, p.~195.
  \href{https://files.ele-math.com/articles/mia-27-15.pdf}{Publisher
  PDF}.
\item
  S. Freewan and M. Hayajneh, \emph{Norm inequalities involving
  geometric means}, arXiv:2401.00337v1, Conjecture 1.4 and the main
  theorem. \href{https://arxiv.org/abs/2401.00337}{Preprint}.
\end{enumerate}

Status check (2026-09-10): the current arXiv record lists v1, December
30, 2023. Its theorem treats \(t=1/2\), \(s\ge2\), \(r\ge1\), \(p>0\),
\(rp\ge1\); the abstract explicitly identifies this as a partial answer
to the weighted conjecture. The original publisher text and the later
preprint statement were compared. Searches for both exact titles, the
authors and ``weighted geometric mean conjecture'', including 2025/2026,
found no full proof or counterexample. Openness is subject to this
search limit.
\end{document}
